\chapter{논의}
\label{ch:discussion}

\ifdraftmode
\begin{center}
\fcolorbox{red}{red!5}{%
\begin{minipage}{0.93\linewidth}
\small
\textcolor{red}{\textbf{본 장의 상태}}\par\smallskip
결과가 존재하지 않으므로 본 장은 결과 \emph{해석}을 담을 수 없다. 대신 결과를
보기 전에 확정해 두는 \textbf{해석 규칙}을 기술한다. 어떤 결과가 나오면 어떤 결론을
내릴 것인지를 사전에 고정하는 것은, 결과를 본 뒤에 유리한 서사를 고르는 것을
막는 유일한 방법이다.
\end{minipage}}
\end{center}
\medskip
\fi

\section{사전에 확정한 해석 규칙}
\label{sec:interpretation-rules}

본 연구의 결과는 네 가지 중 하나로 나올 수 있으며, 각 경우의 결론을 미리 정한다.

\paragraph{경우 1: 빠르고 품질 유지.}
식 \eqref{eq:cost-hypothesis}이 성립하고, 저장량이 $\CR>1$이며, 품질이 independent
2DGS 대비 사전 마진 안에 들어오는 경우. 이는 계산 재사용과 compact representation
가설을 지지한다. 다만 이것만으로 mesh보다 낫다는 결론은 나오지 않는다. RQ5가 별도로
답해야 한다.

\paragraph{경우 2: 빠르지만 품질 저하.}
비용은 줄지만 품질이 마진을 벗어나는 경우. 결론은 ``방법이 실패했다''가 아니라
``절충점이 어디인가''이다. 잔차 계수, surfel 수, 정사영 보정 횟수에 따른 Pareto
curve를 제시하여 어느 계산 예산에서 어떤 표현이 유리한지 보인다.

\paragraph{경우 3: 2DGS가 mesh보다 이득 없음.}
RQ5의 어떤 지표에서도 2DGS가 개선하지 못하는 경우. 이 경우의 결론은 명확하다.
\emph{단색 또는 단순 표면 시각화에는 mesh가 더 적절하다}. 이는 유효한 연구 결과이며
본 논문은 이를 부정적 결과로 취급하지 않는다. 그 경우에도 제 \ref{ch:method},
\ref{ch:error}장의 이론적 기여는 유지된다. 법선 정사영과 오차 분석은 primitive 선택과
독립적이며, mesh 정점에도 같은 방식으로 적용된다.

\paragraph{경우 4: 정사영이 병목.}
$C_{\mathrm{project}}$가 지배적인 경우. surfel 수, 정사영 반복 수, 표면 자료구조에
따른 손익분기점을 분석한다. 특히 명제 \ref{prop:newton-steps}이 예측하는 $K_p\in\{1,2\}$가
실제로 충분한지가 이 경우의 핵심이다.

어떤 결과가 나오더라도 ``표면 변화만으로 어느 수준의 동적 Gaussian 표현을 저비용으로
재사용할 수 있는지''는 밝힐 수 있다.

\section{결과를 보기 전에 정해야 하는 것}

\caveat{%
제 \ref{sec:statistics}절에서 밝힌 대로 비열등성 마진은 본 논문에서 정하지 않았다.
5명 규모 pilot의 측정 noise를 보고 고정해야 하며, 결과를 본 뒤에 정하면 검정이
무의미해진다. 따라서 경우 1과 경우 2를 구분하는 기준선 자체가 아직 존재하지 않는다.
이는 원고의 미완성이 아니라 실험 설계의 필수 순서이다.}

\section{예상되는 실패 모드}

이론이 미리 알려 주는 취약점이 있으며, 실험은 이들을 \emph{확인하러} 설계되었다.

\begin{itemize}[leftmargin=2em]
  \item \textbf{큰 프레임 간 이동.} 법선 정사영은 인접 프레임 사이의 이동이 작고 동일한
  표면 분기가 유지될 때 가장 안정적이다(제 \ref{sec:tangential}절). 수축 프로파일이
  가파른 구간에서 trust region 발동 빈도가 오를 것으로 예상된다.
  \item \textbf{높은 곡률.} 명제 \ref{prop:planar-disk}의 $O(\kappa s^2)$ 오차는 심첨부와
  papillary muscle 근방에서 가장 크다. 곡률 적응 scale(식 \eqref{eq:curvature-scale})이
  이를 억제하는지가 v3 단계의 핵심 검증이다.
  \item \textbf{밝기 비균질.} papillary muscle과 blood pool이 섞이면 Chan--Vese의
  조각마다 상수 가정이 깨진다. 팬텀은 이를 의도적으로 포함하므로 분할 정확도의 저하가
  관측될 것으로 예상된다.
  \item \textbf{방향 표류.} 명제 \ref{prop:propagation}에 의해 위치는 재설정되지만 방향은
  누적된다. 긴 시퀀스에서 $\Eflicker$가 상승하는지, 그리고 등방 원반이 실제로 이를
  무해하게 만드는지(명제 \ref{prop:gauge}) 확인해야 한다.
  \item \textbf{임의 탐색.} 제 \ref{sec:seek-problem}절의 순차 대 정준 정사영 격차.
  이것이 크면 식 \eqref{eq:store}의 저장량 주장과 임의 탐색 기능이 양립하지 않는다.
\end{itemize}

\section{이론과 구현의 관계에 대하여}

본 연구에서 구현은 이론의 부속물이 아니었다. 구현 과정에서 이론 문서가 다루지 않은
세 가지 질문이 강제로 드러났다.

첫째, \emph{감독 신호가 정확히 무엇인가}. 제안서는 ``알려진 slice plane의 mask, depth,
normal, intensity''로 제한했으나 그것을 계산 가능한 목표로 바꾸는 정의가 없었다.
제 \ref{sec:impl-decisions}절의 광선 행진 정의가 그 공백을 채운다.

둘째, \emph{기하가 최적화되지 않음을 어떻게 보장하는가}. 이론적으로는 자명하지만
구현에서는 설정 하나로 깨질 수 있다. 버퍼 등록이 이를 구조적으로 만든다.

셋째, \emph{저장 표현이 임의 탐색을 지원하는가}. 식 \eqref{eq:store}만 보면 자명해
보이지만, 전처리의 순차 정사영과 재생의 직접 정사영이 다르다는 사실은 구현하지
않으면 드러나지 않는다.

이 세 가지는 이론의 오류가 아니라 이론이 명시하지 않은 자유도였다. 구현이 그것을
특정 방식으로 고정했고, 본 논문은 그 선택을 명시적으로 기록하였다.
